\documentclass[14 pt]{extarticle}

	\usepackage[frenchb]{babel}
	\usepackage[utf8]{inputenc}  
	\usepackage[T1]{fontenc}
	\usepackage{amssymb}
	\usepackage[mathscr]{euscript}
	\usepackage{stmaryrd}
	\usepackage{amsmath}
	\usepackage{tikz}
	\usepackage[all,cmtip]{xy}
	\usepackage{amsthm}
	\usepackage{varioref}
	\usepackage{geometry}
	\geometry{a4paper}
	\usepackage{lmodern}
	\usepackage{hyperref}
	\usepackage{array}
	 \usepackage{fancyhdr}
	 \usepackage{float}
\renewcommand{\theenumi}{\alph{enumi})}
	\pagestyle{fancy}
	\theoremstyle{plain}
	\fancyfoot[C]{} 
	\fancyhead[L]{Contrôle}
	\fancyhead[R]{2025-2026}\geometry{
 a4paper,
 total={170mm,257mm},
 left=20mm,
 top=20mm,
 }
	
	
	\title{Contrôle chapitre 3}
	\date{}
	\begin{document}

\begin{center}{\Large Contrôle chapitre 3}\\ 
\emph{Toutes les réponses seront rédigées sur votre copie soigneusement présentée. Calculatrice interdite.}
 \end{center}
 \subsection*{Exercice 1 (4 points)}
 Réduire au même dénominateur : 
 
 \begin{enumerate}
 \item $\frac12$ et $\frac54$
 \item $\frac13$ et $\frac17$
 \item $\frac3{10}$ et $\frac2{15}$
 \item $\frac7{40}$, $\frac3{20}$ et $\frac{11}{16}$
 \end{enumerate}

\subsection*{Exercice 2 (4 points)}
Réduire les fractions suivantes au même dénominateur, et les classer dans l'ordre croissant. 
\[ \frac4{15}, \frac5{12}, \frac56, \frac7{10}\]
\subsection*{Exercice 3 (4 points)}

Tracer un axe gradué de $19$ carreaux avec $1$ au sixième carreau et y placer les fractions suivantes : 
\[ \frac56, \frac{21}{18}, \frac{81}{54}, \frac{108}{36}, \frac5{10}\]

\subsection*{Exercice 4 (4 points)}

Pour chacune des fractions suivantes, donner son écriture décimale et l'exprimer comme un pourcentage, arrondi à l'unité. 
\begin{enumerate}
\item $\frac43$,
\item $\frac25$,
\item $\frac8{11}$,
\item $\frac{2}{15}$. 
\end{enumerate}

\subsection*{Exercice 5 (4 points)}
Regrouper les fractions suivantes par fractions égales. 

\[\frac34 ,  \frac{60}{72},  \frac{12}{27}
,\frac56,
 \frac{15}{20},
 \frac{16}{36},
\frac58, \frac68 
 \]
\newpage


\begin{center}{\Large Contrôle chapitre 3}\\ 
\emph{Toutes les réponses seront rédigées sur votre copie soigneusement présentée. Calculatrice interdite.}
 \end{center}
 \subsection*{Exercice 1 (4 points)}
 Réduire au même dénominateur : 
 
 \begin{enumerate}
 \item $\frac32$ et $\frac54$
 \item $\frac43$ et $\frac18$
 \item $\frac2{14}$ et $\frac4{21}$
 \item $\frac7{40}$, $\frac3{20}$ et $\frac{11}{16}$
 \end{enumerate}

\subsection*{Exercice 2 (4 points)}
Réduire les fractions suivantes au même dénominateur, et les classer dans l'ordre croissant. 
\[ \frac3{15}, \frac5{12}, \frac56, \frac7{10}\]
\subsection*{Exercice 3 (4 points)}

Tracer un axe gradué de $19$ carreaux avec $1$ au sixième carreau et y placer les fractions suivantes : 
\[ \frac76, \frac{12}{18}, \frac{45}{54}, \frac{72}{36}, \frac{10}{15}\]

\subsection*{Exercice 4 (4 points)}

Pour chacune des fractions suivantes, donner son écriture décimale et l'exprimer comme un pourcentage, arrondi à l'unité. 
\begin{enumerate}
\item $\frac23$,
\item $\frac65$,
\item $\frac2{11}$,
\item $\frac{7}{15}$. 
\end{enumerate}

\subsection*{Exercice 5 (4 points)}
Regrouper les fractions suivantes par fractions égales. 

\[\frac34 ,  \frac{30}{36},  \frac{24}{54}
,\frac56,
 \frac{15}{20},
 \frac{16}{36},
\frac58, \frac68 
 \]
 	\end{document}
